\documentclass{beamer}
% =====================================================================
% finance-beamer.tex — shared Beamer style for the LLM-in-Finance course
% =====================================================================
% Reusable slide style. Each deck \input's this immediately after
% \documentclass{beamer}. It provides:
%   * the Madrid theme + book colour palette (matches the printed book)
%   * two book-style framing blocks:
%       \begin{bigpicture} ... \end{bigpicture}   ("the bigger picture")
%       \begin{underhood}   ... \end{underhood}    ("under the hood")
%     These mirror the book's `context` and `deepdive` tcolorboxes so the
%     same intuition-vs-mechanism scaffold the reader meets in the book
%     reappears on the slides.
%   * a Python listings style for code frames
%   * appendix conventions: \appendixstart drops a divider and switches
%     to non-numbered backup frames so "extra / less central" material
%     lives after the main talk without inflating the slide count.
%
% Usage:
%   \documentclass{beamer}
%   % =====================================================================
% finance-beamer.tex — shared Beamer style for the LLM-in-Finance course
% =====================================================================
% Reusable slide style. Each deck \input's this immediately after
% \documentclass{beamer}. It provides:
%   * the Madrid theme + book colour palette (matches the printed book)
%   * two book-style framing blocks:
%       \begin{bigpicture} ... \end{bigpicture}   ("the bigger picture")
%       \begin{underhood}   ... \end{underhood}    ("under the hood")
%     These mirror the book's `context` and `deepdive` tcolorboxes so the
%     same intuition-vs-mechanism scaffold the reader meets in the book
%     reappears on the slides.
%   * a Python listings style for code frames
%   * appendix conventions: \appendixstart drops a divider and switches
%     to non-numbered backup frames so "extra / less central" material
%     lives after the main talk without inflating the slide count.
%
% Usage:
%   \documentclass{beamer}
%   % =====================================================================
% finance-beamer.tex — shared Beamer style for the LLM-in-Finance course
% =====================================================================
% Reusable slide style. Each deck \input's this immediately after
% \documentclass{beamer}. It provides:
%   * the Madrid theme + book colour palette (matches the printed book)
%   * two book-style framing blocks:
%       \begin{bigpicture} ... \end{bigpicture}   ("the bigger picture")
%       \begin{underhood}   ... \end{underhood}    ("under the hood")
%     These mirror the book's `context` and `deepdive` tcolorboxes so the
%     same intuition-vs-mechanism scaffold the reader meets in the book
%     reappears on the slides.
%   * a Python listings style for code frames
%   * appendix conventions: \appendixstart drops a divider and switches
%     to non-numbered backup frames so "extra / less central" material
%     lives after the main talk without inflating the slide count.
%
% Usage:
%   \documentclass{beamer}
%   \input{../_style/finance-beamer.tex}
%   \title{...}\subtitle{...}\author{...}\institute{...}\date{\today}
%   \begin{document} ... \end{document}
% =====================================================================

\usetheme{Madrid}
\usecolortheme{whale}
\setbeamertemplate{navigation symbols}{}
\setbeamertemplate{caption}[numbered]
\setbeamercovered{transparent}

% --- Density controls: keep dense, equation-heavy frames on one slide ---
% Slightly smaller body text + tighter lists/blocks. Frame titles and the
% Madrid chrome keep their normal size, so the deck still reads as Madrid,
% just with more vertical room for content.
\setbeamerfont{normal text}{size=\small}
\setbeamertemplate{itemize/enumerate body begin}{\small}
\setbeamertemplate{itemize/enumerate subbody begin}{\footnotesize}
% Tighter block spacing.
\setbeamertemplate{block begin}{%
  \par\vskip2pt%
  \begin{beamercolorbox}[colsep*=.75ex,leftskip=1ex,rightskip=1ex]{block title}%
    \usebeamerfont*{block title}\insertblocktitle%
  \end{beamercolorbox}%
  {\parskip0pt\vskip-1pt}%
  \begin{beamercolorbox}[colsep*=.75ex,vmode,leftskip=1ex,rightskip=1ex]{block body}%
  \vskip1pt}
\setbeamertemplate{block end}{\end{beamercolorbox}\vskip2pt}

\usepackage{amsmath, amssymb}
\usepackage{booktabs}
\usepackage{xcolor}
\usepackage{listings}
\usepackage[skins, breakable]{tcolorbox}

% --- Book colour palette (kept in sync with book/preamble.tex) ---
\definecolor{linkblue}{RGB}{14, 75, 140}
\definecolor{deepdivebg}{RGB}{235, 244, 255}
\definecolor{narrativebg}{RGB}{255, 252, 240}
\definecolor{narrativerule}{RGB}{180, 130, 40}
\definecolor{steelblue}{RGB}{70, 130, 180}
\definecolor{forestgreen}{RGB}{34, 139, 34}
\definecolor{crimson}{RGB}{220, 20, 60}
\definecolor{darkorange}{RGB}{255, 140, 0}
\definecolor{codebg}{RGB}{248, 248, 248}
\definecolor{coderule}{RGB}{210, 210, 210}
\definecolor{keywordblue}{RGB}{0, 100, 180}
\definecolor{commentgray}{RGB}{120, 120, 120}
\definecolor{stringgreen}{RGB}{30, 130, 30}

% Tie the Madrid structural colour to the book's link blue.
\setbeamercolor{structure}{fg=linkblue}
\setbeamercolor{title}{fg=white}
\setbeamercolor{frametitle}{fg=white}

% --- "the bigger picture" : narrative / intuition framing -----------
\newtcolorbox{bigpicture}{%
  enhanced, breakable, boxsep=2pt,
  colback  = narrativebg,
  colframe = narrativerule,
  boxrule  = 0pt, toprule = 1.4pt, bottomrule = 0.4pt,
  leftrule = 0pt, rightrule = 0pt, arc = 0pt,
  left = 6pt, right = 6pt, top = 4pt, bottom = 5pt,
  fonttitle = \footnotesize\scshape\color{narrativerule},
  title = {the bigger picture}, attach title to upper = {\par\smallskip},
}

% --- "under the hood" : the mechanism / technical detail ------------
\newtcolorbox{underhood}{%
  enhanced, breakable, boxsep=2pt,
  colback  = deepdivebg,
  colframe = linkblue,
  boxrule  = 0pt, toprule = 1.4pt, bottomrule = 0.4pt,
  leftrule = 0pt, rightrule = 0pt, arc = 0pt,
  left = 6pt, right = 6pt, top = 4pt, bottom = 5pt,
  fonttitle = \footnotesize\scshape\color{linkblue},
  title = {under the hood}, attach title to upper = {\par\smallskip},
}

% --- Python code style ---------------------------------------------
\lstset{
  basicstyle    = \ttfamily\footnotesize,
  keywordstyle  = \color{keywordblue}\bfseries,
  commentstyle  = \color{commentgray}\itshape,
  stringstyle   = \color{stringgreen},
  showstringspaces = false,
  language      = Python,
  breaklines    = true,
  backgroundcolor = \color{codebg},
  frame         = single,
  rulecolor     = \color{coderule},
  numbers       = none,
  aboveskip     = 4pt, belowskip = 4pt,
}

% --- Appendix conventions ------------------------------------------
% \appendixstart{Title} : begins the "extra material" appendix. Frames
% after it are not counted toward the main deck's frame total, keeping
% the core talk lean while preserving the deeper / optional content.
\newcommand{\appendixstart}[1]{%
  \appendix
  \setbeamertemplate{frametitle continuation}{}%
  \begin{frame}[noframenumbering]
    \begin{center}
      {\Large\scshape\color{linkblue} Appendix}\\[4pt]
      {\large #1}\\[10pt]
      {\footnotesize Extra / optional material — beyond the core lecture.}
    \end{center}
  \end{frame}
}
% Use \begin{frame}[noframenumbering]{...} for each appendix frame.

%   \title{...}\subtitle{...}\author{...}\institute{...}\date{\today}
%   \begin{document} ... \end{document}
% =====================================================================

\usetheme{Madrid}
\usecolortheme{whale}
\setbeamertemplate{navigation symbols}{}
\setbeamertemplate{caption}[numbered]
\setbeamercovered{transparent}

% --- Density controls: keep dense, equation-heavy frames on one slide ---
% Slightly smaller body text + tighter lists/blocks. Frame titles and the
% Madrid chrome keep their normal size, so the deck still reads as Madrid,
% just with more vertical room for content.
\setbeamerfont{normal text}{size=\small}
\setbeamertemplate{itemize/enumerate body begin}{\small}
\setbeamertemplate{itemize/enumerate subbody begin}{\footnotesize}
% Tighter block spacing.
\setbeamertemplate{block begin}{%
  \par\vskip2pt%
  \begin{beamercolorbox}[colsep*=.75ex,leftskip=1ex,rightskip=1ex]{block title}%
    \usebeamerfont*{block title}\insertblocktitle%
  \end{beamercolorbox}%
  {\parskip0pt\vskip-1pt}%
  \begin{beamercolorbox}[colsep*=.75ex,vmode,leftskip=1ex,rightskip=1ex]{block body}%
  \vskip1pt}
\setbeamertemplate{block end}{\end{beamercolorbox}\vskip2pt}

\usepackage{amsmath, amssymb}
\usepackage{booktabs}
\usepackage{xcolor}
\usepackage{listings}
\usepackage[skins, breakable]{tcolorbox}

% --- Book colour palette (kept in sync with book/preamble.tex) ---
\definecolor{linkblue}{RGB}{14, 75, 140}
\definecolor{deepdivebg}{RGB}{235, 244, 255}
\definecolor{narrativebg}{RGB}{255, 252, 240}
\definecolor{narrativerule}{RGB}{180, 130, 40}
\definecolor{steelblue}{RGB}{70, 130, 180}
\definecolor{forestgreen}{RGB}{34, 139, 34}
\definecolor{crimson}{RGB}{220, 20, 60}
\definecolor{darkorange}{RGB}{255, 140, 0}
\definecolor{codebg}{RGB}{248, 248, 248}
\definecolor{coderule}{RGB}{210, 210, 210}
\definecolor{keywordblue}{RGB}{0, 100, 180}
\definecolor{commentgray}{RGB}{120, 120, 120}
\definecolor{stringgreen}{RGB}{30, 130, 30}

% Tie the Madrid structural colour to the book's link blue.
\setbeamercolor{structure}{fg=linkblue}
\setbeamercolor{title}{fg=white}
\setbeamercolor{frametitle}{fg=white}

% --- "the bigger picture" : narrative / intuition framing -----------
\newtcolorbox{bigpicture}{%
  enhanced, breakable, boxsep=2pt,
  colback  = narrativebg,
  colframe = narrativerule,
  boxrule  = 0pt, toprule = 1.4pt, bottomrule = 0.4pt,
  leftrule = 0pt, rightrule = 0pt, arc = 0pt,
  left = 6pt, right = 6pt, top = 4pt, bottom = 5pt,
  fonttitle = \footnotesize\scshape\color{narrativerule},
  title = {the bigger picture}, attach title to upper = {\par\smallskip},
}

% --- "under the hood" : the mechanism / technical detail ------------
\newtcolorbox{underhood}{%
  enhanced, breakable, boxsep=2pt,
  colback  = deepdivebg,
  colframe = linkblue,
  boxrule  = 0pt, toprule = 1.4pt, bottomrule = 0.4pt,
  leftrule = 0pt, rightrule = 0pt, arc = 0pt,
  left = 6pt, right = 6pt, top = 4pt, bottom = 5pt,
  fonttitle = \footnotesize\scshape\color{linkblue},
  title = {under the hood}, attach title to upper = {\par\smallskip},
}

% --- Python code style ---------------------------------------------
\lstset{
  basicstyle    = \ttfamily\footnotesize,
  keywordstyle  = \color{keywordblue}\bfseries,
  commentstyle  = \color{commentgray}\itshape,
  stringstyle   = \color{stringgreen},
  showstringspaces = false,
  language      = Python,
  breaklines    = true,
  backgroundcolor = \color{codebg},
  frame         = single,
  rulecolor     = \color{coderule},
  numbers       = none,
  aboveskip     = 4pt, belowskip = 4pt,
}

% --- Appendix conventions ------------------------------------------
% \appendixstart{Title} : begins the "extra material" appendix. Frames
% after it are not counted toward the main deck's frame total, keeping
% the core talk lean while preserving the deeper / optional content.
\newcommand{\appendixstart}[1]{%
  \appendix
  \setbeamertemplate{frametitle continuation}{}%
  \begin{frame}[noframenumbering]
    \begin{center}
      {\Large\scshape\color{linkblue} Appendix}\\[4pt]
      {\large #1}\\[10pt]
      {\footnotesize Extra / optional material — beyond the core lecture.}
    \end{center}
  \end{frame}
}
% Use \begin{frame}[noframenumbering]{...} for each appendix frame.

%   \title{...}\subtitle{...}\author{...}\institute{...}\date{\today}
%   \begin{document} ... \end{document}
% =====================================================================

\usetheme{Madrid}
\usecolortheme{whale}
\setbeamertemplate{navigation symbols}{}
\setbeamertemplate{caption}[numbered]
\setbeamercovered{transparent}

% --- Density controls: keep dense, equation-heavy frames on one slide ---
% Slightly smaller body text + tighter lists/blocks. Frame titles and the
% Madrid chrome keep their normal size, so the deck still reads as Madrid,
% just with more vertical room for content.
\setbeamerfont{normal text}{size=\small}
\setbeamertemplate{itemize/enumerate body begin}{\small}
\setbeamertemplate{itemize/enumerate subbody begin}{\footnotesize}
% Tighter block spacing.
\setbeamertemplate{block begin}{%
  \par\vskip2pt%
  \begin{beamercolorbox}[colsep*=.75ex,leftskip=1ex,rightskip=1ex]{block title}%
    \usebeamerfont*{block title}\insertblocktitle%
  \end{beamercolorbox}%
  {\parskip0pt\vskip-1pt}%
  \begin{beamercolorbox}[colsep*=.75ex,vmode,leftskip=1ex,rightskip=1ex]{block body}%
  \vskip1pt}
\setbeamertemplate{block end}{\end{beamercolorbox}\vskip2pt}

\usepackage{amsmath, amssymb}
\usepackage{booktabs}
\usepackage{xcolor}
\usepackage{listings}
\usepackage[skins, breakable]{tcolorbox}

% --- Book colour palette (kept in sync with book/preamble.tex) ---
\definecolor{linkblue}{RGB}{14, 75, 140}
\definecolor{deepdivebg}{RGB}{235, 244, 255}
\definecolor{narrativebg}{RGB}{255, 252, 240}
\definecolor{narrativerule}{RGB}{180, 130, 40}
\definecolor{steelblue}{RGB}{70, 130, 180}
\definecolor{forestgreen}{RGB}{34, 139, 34}
\definecolor{crimson}{RGB}{220, 20, 60}
\definecolor{darkorange}{RGB}{255, 140, 0}
\definecolor{codebg}{RGB}{248, 248, 248}
\definecolor{coderule}{RGB}{210, 210, 210}
\definecolor{keywordblue}{RGB}{0, 100, 180}
\definecolor{commentgray}{RGB}{120, 120, 120}
\definecolor{stringgreen}{RGB}{30, 130, 30}

% Tie the Madrid structural colour to the book's link blue.
\setbeamercolor{structure}{fg=linkblue}
\setbeamercolor{title}{fg=white}
\setbeamercolor{frametitle}{fg=white}

% --- "the bigger picture" : narrative / intuition framing -----------
\newtcolorbox{bigpicture}{%
  enhanced, breakable, boxsep=2pt,
  colback  = narrativebg,
  colframe = narrativerule,
  boxrule  = 0pt, toprule = 1.4pt, bottomrule = 0.4pt,
  leftrule = 0pt, rightrule = 0pt, arc = 0pt,
  left = 6pt, right = 6pt, top = 4pt, bottom = 5pt,
  fonttitle = \footnotesize\scshape\color{narrativerule},
  title = {the bigger picture}, attach title to upper = {\par\smallskip},
}

% --- "under the hood" : the mechanism / technical detail ------------
\newtcolorbox{underhood}{%
  enhanced, breakable, boxsep=2pt,
  colback  = deepdivebg,
  colframe = linkblue,
  boxrule  = 0pt, toprule = 1.4pt, bottomrule = 0.4pt,
  leftrule = 0pt, rightrule = 0pt, arc = 0pt,
  left = 6pt, right = 6pt, top = 4pt, bottom = 5pt,
  fonttitle = \footnotesize\scshape\color{linkblue},
  title = {under the hood}, attach title to upper = {\par\smallskip},
}

% --- Python code style ---------------------------------------------
\lstset{
  basicstyle    = \ttfamily\footnotesize,
  keywordstyle  = \color{keywordblue}\bfseries,
  commentstyle  = \color{commentgray}\itshape,
  stringstyle   = \color{stringgreen},
  showstringspaces = false,
  language      = Python,
  breaklines    = true,
  backgroundcolor = \color{codebg},
  frame         = single,
  rulecolor     = \color{coderule},
  numbers       = none,
  aboveskip     = 4pt, belowskip = 4pt,
}

% --- Appendix conventions ------------------------------------------
% \appendixstart{Title} : begins the "extra material" appendix. Frames
% after it are not counted toward the main deck's frame total, keeping
% the core talk lean while preserving the deeper / optional content.
\newcommand{\appendixstart}[1]{%
  \appendix
  \setbeamertemplate{frametitle continuation}{}%
  \begin{frame}[noframenumbering]
    \begin{center}
      {\Large\scshape\color{linkblue} Appendix}\\[4pt]
      {\large #1}\\[10pt]
      {\footnotesize Extra / optional material — beyond the core lecture.}
    \end{center}
  \end{frame}
}
% Use \begin{frame}[noframenumbering]{...} for each appendix frame.


\title{Practical Session 3: LoRA, Scaling, and Fine-Tuning}
\subtitle{Hands-On: Parameter Counting, Chinchilla Budgeting, and FinBERT Sentiment}
\author{Juan F.\ Imbet}
\institute{Paris Dauphine -- PSL University}
\date{Large Language Models in Finance}

\begin{document}

% ============================================================
% FRAME 1: Title
% ============================================================
\begin{frame}
\frametitle{}
\titlepage
\end{frame}

% ============================================================
% FRAME 2: Session overview
% ============================================================
\begin{frame}
\frametitle{Session Overview}

\textbf{Duration:} 1 hour

\medskip
\begin{enumerate}
  \item \textbf{Recap} (5 min) --- LoRA formula, Chinchilla rule, CLM/MLM objectives
  \item \textbf{Problem 1} (15 min) --- LoRA parameter counting on 7B LLaMA
  \item \textbf{Problem 2} (10 min) --- Chinchilla compute budgeting
  \item \textbf{Case Study} (20 min) --- Fine-tune FinBERT sentiment with PEFT/LoRA;
        answer three diagnostic questions
  \item \textbf{Discussion} (5 min) --- When to choose each adaptation strategy?
  \item \textbf{Solutions} (5 min) --- Walkthrough of Problems 1 and 2
\end{enumerate}

\medskip
\begin{block}{What you will practise}
  Translating scaling-law arithmetic into practical hardware decisions; applying
  LoRA configuration choices; running a real fine-tuning loop on a finance task.
\end{block}
\end{frame}

% ============================================================
% FRAME 3: Recap -- LoRA
% ============================================================
\begin{frame}
\frametitle{Recap: LoRA Formula}

For a frozen pre-trained weight matrix $W_0 \in \mathbb{R}^{d \times k}$:
\[
  W = W_0 + BA, \quad B \in \mathbb{R}^{d \times r},\; A \in \mathbb{R}^{r \times k},\;
  r \ll \min(d,k)
\]

\medskip
\textbf{Trainable parameter count per adapted layer:}
\[
  \Delta\theta = r(d + k)
\]

\medskip
\textbf{Key properties:}
\begin{itemize}
  \item $W_0$ is frozen; only $B$ and $A$ receive gradients
  \item Initialise: $A \sim \mathcal{N}(0, \sigma^2)$, $B = 0$ so $BA = 0$ at step 0
  \item At inference: merge $W \leftarrow W_0 + BA$ --- zero overhead vs.\ base model
  \item Rank $r$ is the key hyperparameter; typical values: 4, 8, 16, 32
\end{itemize}

\medskip
\begin{alertblock}{You will need this for Problem 1}
  Summing $r(d+k)$ over all adapted layers gives total trainable parameter count.
\end{alertblock}
\end{frame}

% ============================================================
% FRAME 4: Recap -- Chinchilla + CLM/MLM
% ============================================================
\begin{frame}
\frametitle{Recap: Chinchilla Rule and Pre-training Objectives}

\textbf{Chinchilla compute-optimal allocation:}
\[
  C \approx 6ND, \qquad \frac{D^*}{N^*} \approx 20
\]
Given budget $C$: solve $6N(20N) = 120N^2 = C$, so
\[
  N^* = \sqrt{\frac{C}{120}}, \qquad D^* = 20 N^*
\]

\medskip
\textbf{Pre-training objectives (summary):}

\begin{center}
\begin{tabular}{lll}
\toprule
\textbf{Objective} & \textbf{Architecture} & \textbf{Finance use} \\
\midrule
CLM (causal) & Decoder-only & Generation, chat \\
MLM (masked) & Encoder-only & Classification, NER \\
Span corruption & Encoder-decoder & QA, summarisation \\
\bottomrule
\end{tabular}
\end{center}

\medskip
\begin{alertblock}{You will need these for Problem 2}
  The formula $N^* = \sqrt{C/120}$ is your main tool for Chinchilla budgeting.
\end{alertblock}
\end{frame}

% ============================================================
% FRAME 5: Problem 1 -- Setup
% ============================================================
\begin{frame}
\frametitle{Problem 1: LoRA Parameter Counting --- Setup}

\textbf{Model:} LLaMA-7B with 32 transformer layers, hidden dimension $d = 4096$.

Each layer contains four linear projection matrices:
\begin{itemize}
  \item Query projection: $W_Q \in \mathbb{R}^{4096 \times 4096}$
  \item Key projection: $W_K \in \mathbb{R}^{4096 \times 4096}$
  \item Value projection: $W_V \in \mathbb{R}^{4096 \times 4096}$
  \item Output projection: $W_O \in \mathbb{R}^{4096 \times 4096}$
\end{itemize}

\textbf{Task:} Apply LoRA with rank $r = 4$ to the \emph{Q and V projections only}.

\medskip
\textbf{Fill in the table below for one layer:}

\begin{center}
\begin{tabular}{lcccc}
\toprule
\textbf{Layer} & \textbf{Matrix} & \textbf{$d$} & \textbf{$k$} & \textbf{$r$} \\
\midrule
Transformer $\ell$ & $W_Q$ & \phantom{4096} & \phantom{4096} & \phantom{4} \\
Transformer $\ell$ & $W_V$ & \phantom{4096} & \phantom{4096} & \phantom{4} \\
\bottomrule
\end{tabular}
\end{center}

\medskip
\small \textit{Hint:} For each adapted matrix the trainable param count is $r(d + k)$.
Sum over both matrices and all 32 layers.
\end{frame}

% ============================================================
% FRAME 6: Problem 1 -- Tasks
% ============================================================
\begin{frame}
\frametitle{Problem 1: Tasks}

Using the setup from the previous frame ($r=4$, Q and V only, 32 layers, $d=k=4096$):

\medskip
\textbf{Task A.} Compute the total number of trainable LoRA parameters.
Show the step-by-step arithmetic.

\medskip
\textbf{Task B.} Express this as a fraction of the 7B total parameter count.
Is the fraction above or below 0.1\%?

\medskip
\textbf{Task C.} Suppose you double the rank from $r = 4$ to $r = 8$.
By what factor does the trainable parameter count change?
Is the relationship linear, quadratic, or something else? Justify.

\medskip
\textbf{Task D.} A colleague suggests using adapter modules instead of LoRA
because adapters have been around longer. Give one concrete reason why LoRA
is preferred for \emph{inference-time deployment} in a high-throughput
trading-floor application where latency is measured in milliseconds.

\medskip
\begin{block}{Time: 15 minutes. Work individually or in pairs.}
\end{block}
\end{frame}

% ============================================================
% FRAME 7: Problem 2 -- Chinchilla budgeting
% ============================================================
\begin{frame}
\frametitle{Problem 2: Chinchilla Compute Budgeting}

\textbf{Scenario:} An investment bank's AI team has a compute budget of
$C = 10^{22}$ FLOPs for training a financial language model.

\medskip
Using $C \approx 6ND$, $D^* \approx 20 N^*$, and therefore $N^* = \sqrt{C/120}$:

\medskip
\textbf{Task A.}
\begin{enumerate}
  \item Compute $N^*$ and $D^*$.
  \item The team has assembled a financial text corpus containing 50B tokens.
        Is the team data-constrained or compute-constrained?
        (Hint: compare your $D^*$ to 50B tokens.)
  \item What should they do given this diagnosis?
\end{enumerate}
\end{frame}

% ============================================================
% FRAME 7b: Problem 2 -- Tasks B and C (continued)
% ============================================================
\begin{frame}
\frametitle{Problem 2: Chinchilla Compute Budgeting (continued)}

\textbf{Task B.} The team instead decides to train a 1B-parameter model.
Using the Chinchilla rule, how many tokens should they train on for a compute-optimal run?
What is the implied FLOP cost?

\medskip
\textbf{Task C.} A senior researcher argues ``we should train the biggest model
possible to maximise accuracy''. Under what conditions is this argument valid?
Under what conditions is it wrong?

\medskip
\begin{block}{Time: 10 minutes.}
\end{block}
\end{frame}

% ============================================================
% FRAME 8: Case Study -- Introduction
% ============================================================
\begin{frame}
\frametitle{Case Study: Fine-Tuning FinBERT with LoRA on Financial Sentiment}

\textbf{Goal:} Load a pre-trained FinBERT model and apply LoRA fine-tuning
on a small set of labelled financial sentences using HuggingFace PEFT.

\medskip
\textbf{Stack:}
\begin{itemize}
  \item \texttt{transformers} --- load \texttt{ProsusAI/finbert}
  \item \texttt{peft} --- apply \texttt{LoraConfig}, wrap model with \texttt{get\_peft\_model}
  \item \texttt{torch} --- training loop (no Trainer class for transparency)
\end{itemize}
\end{frame}

% ============================================================
% FRAME 8b: Case Study -- Introduction (continued)
% ============================================================
\begin{frame}
\frametitle{Case Study: FinBERT with LoRA --- What You Will Do}

\textbf{Dataset:} 5 labelled financial sentences (positive/negative/neutral).
This is intentionally tiny --- the goal is to observe the mechanics, not maximise accuracy.

\medskip
\begin{block}{What you will do}
  \begin{enumerate}
    \item Run the code on the next frame
    \item Observe how many parameters are trainable vs.\ frozen
    \item Answer the diagnostic questions on the following frame
  \end{enumerate}
\end{block}

\medskip
\begin{alertblock}{Prerequisites}
  \texttt{pip install transformers peft torch}\\
  No API key required for this case study.
\end{alertblock}
\end{frame}

% ============================================================
% FRAME 9: Case Study -- Code
% ============================================================
\begin{frame}[fragile]
\frametitle{Case Study: LoRA on FinBERT --- Code (Part 1a: Imports + Data)}

{\footnotesize
\begin{lstlisting}
import torch
from transformers import AutoTokenizer, AutoModelForSequenceClassification
from peft import LoraConfig, get_peft_model, TaskType

MODEL = "ProsusAI/finbert"
SENTENCES = [
    ("Revenue increased 15% year-over-year.", 0),        # positive
    ("The firm reported a net loss of $200M.", 1),        # negative
    ("Operating margins remained broadly stable.", 2),    # neutral
    ("Strong demand drove record quarterly earnings.", 0),# positive
    ("Guidance was withdrawn amid macro uncertainty.", 1),# negative
]
\end{lstlisting}
}
\end{frame}

% ============================================================
% FRAME 9 (Part 1b): Case Study -- Code (Load + LoRA config)
% ============================================================
\begin{frame}[fragile]
\frametitle{Case Study: LoRA on FinBERT --- Code (Part 1b: Load + LoRA Config)}

{\footnotesize
\begin{lstlisting}
tokenizer = AutoTokenizer.from_pretrained(MODEL)
model = AutoModelForSequenceClassification.from_pretrained(MODEL, num_labels=3)

lora_cfg = LoraConfig(
    task_type=TaskType.SEQ_CLS,
    r=8,
    lora_alpha=16,
    target_modules=["query", "value"],
    lora_dropout=0.05,
    bias="none",
)
model = get_peft_model(model, lora_cfg)
model.print_trainable_parameters()  # observe trainable % here
\end{lstlisting}
}
\end{frame}

% ============================================================
% FRAME 9b: Case Study -- Code (Part 2: Train + Eval)
% ============================================================
\begin{frame}[fragile]
\frametitle{Case Study: LoRA on FinBERT --- Code (Part 2: Train + Eval)}

{\footnotesize
\begin{lstlisting}
texts, labels = zip(*SENTENCES)
enc = tokenizer(list(texts), padding=True, truncation=True, return_tensors="pt")
label_tensor = torch.tensor(labels)

optimizer = torch.optim.AdamW(model.parameters(), lr=2e-4)
model.train()
for epoch in range(3):
    out = model(**enc, labels=label_tensor)
    out.loss.backward(); optimizer.step(); optimizer.zero_grad()
    print(f"Epoch {epoch+1}: loss={out.loss.item():.4f}")

model.eval()
with torch.no_grad():
    logits = model(**enc).logits
    preds = logits.argmax(-1)
    id2label = {0: "positive", 1: "negative", 2: "neutral"}
    for sent, pred in zip(texts, preds):
        print(f"[{id2label[pred.item()]}] {sent}")
\end{lstlisting}
}
\end{frame}

% ============================================================
% FRAME 10: Case Study -- Diagnostic questions
% ============================================================
\begin{frame}
\frametitle{Case Study: Diagnostic Questions}

Run the code and observe the output, then answer:

\medskip
\textbf{Q1.} \texttt{model.print\_trainable\_parameters()} reports the number of
trainable parameters.  Compare this to the total FinBERT parameter count
($\approx 110$M for BERT-large variants).  How does the trainable fraction
compare to the theoretical value for LoRA $r=8$ applied to query and value
matrices across all 12 BERT layers?

\medskip
\textbf{Q2.} Change \texttt{r=8} to \texttt{r=4} and then to \texttt{r=16}.
Run the three-epoch training loop for each.  Do the final predictions change?
Does loss decrease faster at higher rank?  What does this suggest about the
relationship between rank and expressiveness for this small dataset?

\medskip
\end{frame}

% ============================================================
% FRAME 10b: Case Study -- Diagnostic questions (continued)
% ============================================================
\begin{frame}
\frametitle{Case Study: Diagnostic Questions (continued)}

\textbf{Q3.} Modify \texttt{target\_modules} to include all four projections:
\texttt{["query", "key", "value", "dense"]}.
How does the trainable parameter count change?
Do the predictions improve or change?  Is the additional compute justified?

\medskip
\begin{block}{Optional extension}
  Replace the 5 sentences with 50 examples from the Financial PhraseBank dataset
  (available on HuggingFace as \texttt{takala/financial\_phrasebank}).
  Evaluate accuracy on a held-out 20\% split at $r = 4$, $8$, $16$.
  Plot accuracy vs.\ trainable parameter count.
\end{block}
\end{frame}

% ============================================================
% FRAME 11: Discussion
% ============================================================
\begin{frame}
\frametitle{Discussion: Choosing the Right Adaptation Strategy}

For each scenario below, identify which strategy you would use
(zero-shot / few-shot ICL / LoRA fine-tune / full fine-tune) and why.

\medskip
\begin{enumerate}
  \item A hedge fund needs to classify 10-K risk-factor sections as
        ``material'' or ``immaterial'' for a specific covenant type.
        They have 50 labelled examples and a 2-hour deadline.
        They use GPT-4o via API.

  \item A tier-1 bank wants to build an internal model for earnings-call
        sentiment that runs on-premise (no external API), processes
        200,000 documents per month, and must meet GDPR data-residency requirements.
        They have 20,000 labelled examples.

  \item A fintech startup is building a general-purpose financial chatbot.
        They have no task-specific labelled data but have access to a
        pre-trained instruction-tuned LLM.

  \item A quantitative research team wants to test whether a 70B model
        can replicate the reasoning of a domain expert on novel bond
        structuring questions.  No labelled data exists.
  \end{enumerate}

\medskip
\begin{block}{Frame your answer around:}
  Labelled data volume, compute budget, latency requirements,
  data-privacy constraints, inference cost, and desired explainability.
\end{block}
\end{frame}

% ============================================================
% FRAME 12: Solution -- Problem 1
% ============================================================
\begin{frame}
\frametitle{Solution: Problem 1 --- LoRA Parameter Counting}

\textbf{Filled table (one layer, $r=4$, Q and V, $d=k=4096$):}

\begin{center}
\begin{tabular}{lcccc}
\toprule
\textbf{Layer} & \textbf{Matrix} & \textbf{$d$} & \textbf{$k$} & \textbf{$r$} \\
\midrule
Transformer $\ell$ & $W_Q$ & 4096 & 4096 & 4 \\
Transformer $\ell$ & $W_V$ & 4096 & 4096 & 4 \\
\bottomrule
\end{tabular}
\end{center}

\medskip
\textbf{Task A:}
\[
  32 \times 2 \times r \times (d + k)
  = 32 \times 2 \times 4 \times 8192
  = 2{,}097{,}152 \approx 2.1\text{M params}
\]

\textbf{Task B:} $2.1\text{M} / 7000\text{M} \approx 0.030\%$ --- well below 0.1\%.

\textbf{Task C:} Doubling $r$ from 4 to 8:
\[
  32 \times 2 \times 8 \times 8192 = 4{,}194{,}304 \approx 4.2\text{M}
\]
The count \emph{doubles} (linear in $r$, since $r(d+k)$ is linear in $r$).

\textbf{Task D:} LoRA merges $W_0 + BA$ into a single weight before deployment,
so the merged model is identical in structure to the base model.
Adapters add a sequential bottleneck module that cannot be trivially removed,
increasing inference latency per token --- unacceptable in latency-sensitive
trading applications.
\end{frame}

% ============================================================
% FRAME 13: Solution -- Problem 2 + wrap-up
% ============================================================
\begin{frame}
\frametitle{Solution: Problem 2 --- Chinchilla Budgeting + Wrap-Up}

\textbf{Task A} ($C = 10^{22}$):
\[
  N^* = \sqrt{\frac{C}{120}} = \sqrt{\frac{10^{22}}{120}}
       = \sqrt{8.33 \times 10^{19}} \approx 2.89 \times 10^8 \approx 288\text{M params}
\]
\[
  D^* = 20 \times N^* \approx 20 \times 288\text{M} = 5.76\text{B tokens}
\]
The team has 50B tokens but $D^* \approx 5.76$B.
$50\text{B} \gg 5.76\text{B}$: they are \emph{compute-constrained, not data-constrained}.
They should train a $\approx 288$M-parameter model on 5.76B tokens; the remaining
44B tokens support a larger model if compute permits.
\end{frame}

% ============================================================
% FRAME 13b: Solution -- Problem 2 (Tasks B/C) + wrap-up
% ============================================================
\begin{frame}
\frametitle{Solution: Problem 2 --- Tasks B/C + Wrap-Up}

\textbf{Task B} (1B-parameter model):
$D^* = 20 \times 1\text{B} = 20\text{B tokens}$; $C = 6 \times 1\text{B} \times 20\text{B} = 1.2 \times 10^{23}$ FLOPs.

\medskip
\textbf{Task C:} Maximise-parameter argument is valid only when inference cost is
negligible and compute is unlimited.  In practice, the 50B-parameter model costs
more to serve at inference time than a Chinchilla-optimal 288M model, and may
\emph{score worse} on held-out loss if under-trained.

\medskip
\begin{block}{Session Summary}
  LoRA param count is linear in $r$; merge eliminates inference overhead.
  Chinchilla: $D^*/N^* \approx 20$; most realistic settings are compute-constrained.
  Next practical: fine-tuning a generative model for financial QA with LoRA.
\end{block}
\end{frame}

% ============================================================
% APPENDIX
% ============================================================
\appendixstart{Stretch Problem: DPO on a Preference Dataset}

\begin{frame}[noframenumbering]{Stretch Problem: From LoRA Fine-Tuning to DPO}

\textbf{Setup.} You have a LoRA-fine-tuned FinBERT-style sentiment model from the
case study. Now consider aligning a small \emph{generative} finance model with
Direct Preference Optimisation (DPO) on preference pairs $(x, y_w, y_l)$.

\medskip
\textbf{Task A.} Recall the DPO objective:
\[
  \mathcal{L}_{\text{DPO}}(\theta) = -\mathbb{E}_{(x,y_w,y_l)}\!\left[
    \log \sigma\!\left(
      \beta \log \tfrac{\pi_\theta(y_w|x)}{\pi_{\text{ref}}(y_w|x)}
      - \beta \log \tfrac{\pi_\theta(y_l|x)}{\pi_{\text{ref}}(y_l|x)}
    \right)\right].
\]
Explain, in one sentence, what raising $\beta$ does to how far $\pi_\theta$ may drift
from $\pi_{\text{ref}}$.

\medskip
\textbf{Task B.} Why does DPO need a frozen reference policy $\pi_{\text{ref}}$ at all?
What failure occurs if you drop it?

\medskip
\textbf{Task C.} Construct three plausible $(y_w \succ y_l)$ preference pairs for a
financial-advice prompt where ``plausible vs.\ accurate'' is the distinguishing axis.

\medskip
\begin{block}{Time: 15 minutes (optional / take-home).}
\end{block}
\end{frame}

\end{document}
